\section{Related Work}

\label{sec:related}
%----------------------------------------------------------------------

\subsection{DeSci Projects}

\textbf{VitaDAO} funds longevity research through token-governed grants. Unlike \Zoo{}, VitaDAO is domain-specific (longevity) and does not provide research infrastructure (computation, data hosting, reproducibility). \textbf{ResearchHub} incentivizes open-access publishing with a token reward for peer review. It focuses on the publication layer, not the research process. \textbf{DeSci Labs} builds tools for research data management but without on-chain provenance or verifiable computation.

\Zoo{} differs from all existing DeSci projects in providing a \emph{complete} research infrastructure stack: from data contribution through computation, review, funding, publication, and formal verification.

\subsection{Reproducibility Platforms}

\textbf{Code Ocean} provides containerized compute capsules for reproducible research. \textbf{Binder} offers free ephemeral environments for Jupyter notebooks. \textbf{Whole Tale} integrates data, code, and compute for reproducibility.

These platforms provide reproducibility but lack: on-chain provenance, decentralized governance, economic incentives for contribution, and formal verification integration.

\subsection{Open Data Initiatives}

\textbf{Zenodo} (CERN) provides open-access data hosting with DOIs. \textbf{Dataverse} offers institutional data repositories. \textbf{GBIF} aggregates biodiversity occurrence records.

These initiatives provide data hosting but not: usage-based compensation, decentralized governance, or AI-assisted analysis.

%----------------------------------------------------------------------
